\documentclass[11pt,a4paper]{article}
\usepackage[margin=1in]{geometry}
\usepackage{times}
\usepackage{latexsym}
\usepackage{graphicx}
\usepackage{amsmath}
\usepackage{amssymb}
\usepackage{booktabs}
\usepackage{multirow}
\usepackage{xcolor}
\usepackage{tabularx}
\usepackage{array}
\usepackage[T1]{fontenc}
\usepackage[utf8]{inputenc}
\usepackage{microtype}
\usepackage[colorlinks=true,citecolor=blue,linkcolor=blue,urlcolor=blue]{hyperref}
\usepackage{natbib}
\renewcommand{\bibsection}{\section*{References}}

\definecolor{siggreen}{rgb}{0.0,0.5,0.0}
\definecolor{trendyellow}{rgb}{0.7,0.5,0.0}
\newcommand{\sig}[1]{\textcolor{siggreen}{\textbf{#1}}}
\newcommand{\trend}[1]{\textcolor{trendyellow}{\textit{#1}}}
\newcommand{\ns}[1]{\textcolor{gray}{#1}}

\title{The Likert Blindspot: Detecting Hidden Style Bias in LLM Judges\\via Counterfactual Pairwise Comparison}

\author{Anonymous}

\date{}

\begin{document}
\maketitle

\begin{abstract}
LLM-as-judge systems are increasingly central to NLP evaluation and alignment training, yet standard Likert pointwise scoring systematically fails to detect style bias.
We introduce a counterfactual pairwise comparison protocol that generates style-varied rewrites of evaluation responses, verifies content preservation via NLI, and measures judge preference through position-controlled forced choice.
Applying this protocol to formality bias across three judges (GPT-4o, Claude-Sonnet, Gemini-Flash) and three benchmarks (LLMBar, AlpacaEval, MT-Bench), we find that Likert scoring reports no bias in all nine conditions (0/9 significant), while pairwise comparison detects significant formality preference in two conditions (GPT-4o$\times$LLMBar: 69.1\%, $p{=}0.006$; Gemini-Flash$\times$LLMBar: 63.8\%, $p{=}0.041$) with supporting trends on AlpacaEval.
A scaled replication ($n{=}500$) confirms this \textit{Likert Blindspot} is primarily a measurement efficiency gap: pairwise comparison requires ${\sim}100$ samples for 80\% power, while binarized Likert requires ${\sim}680$ --- a ${\sim}7\times$ sample efficiency advantage.
We provide power analysis tables, rewriter robustness checks, and practical guidance for practitioners auditing LLM judge style bias.
\end{abstract}

%------------------------------------------------------------------
\section{Introduction}
\label{sec:intro}
%------------------------------------------------------------------

LLM-as-judge evaluation has become a core component of the NLP evaluation stack, powering leaderboards such as Chatbot Arena \cite{zheng2023judging}, training data curation pipelines, and alignment reward signals.
A growing body of evidence shows that LLM judges exhibit systematic \emph{style bias}: preference for responses with particular surface characteristics---verbosity, formality, markdown formatting---independent of content quality \cite{sosbench2025,casse2026,soumik2026}.
The correlation between response style and judge scores can be remarkably strong ($R{>}0.99$ for length in SOS-Bench), raising concerns that downstream systems optimize for style rather than substance.

Standard evaluation practice relies on Likert pointwise scoring, where a judge assigns an integer quality rating to each response independently.
We find that this default measurement paradigm has a systematic blind spot for style bias.
On the same counterfactual data, Likert scoring reports no formality bias ($p{=}0.332$), while position-controlled pairwise comparison detects a significant formal-style preference (69.1\% win rate, $p{=}0.006$) for GPT-4o on LLMBar.
This \emph{Likert Blindspot} is not a single-condition anomaly: across nine judge$\times$benchmark combinations, continuous Likert ATE is non-significant in every case (0/9), whereas pairwise comparison reveals significant or trending effects in four conditions.

Our contributions are:
\begin{enumerate}
    \item A counterfactual pairwise comparison protocol combining style-targeted rewriting, NLI-based content verification, and position-controlled forced choice for measuring black-box LLM judge style bias (\S\ref{sec:method}).
    \item Systematic quantification of the Likert Blindspot across three judges and three benchmarks, demonstrating that pairwise comparison provides a ${\sim}1.7\times$ effect amplification and ${\sim}7\times$ sample efficiency advantage over Likert-based approaches (\S\ref{sec:exp_sensitivity}, \S\ref{sec:exp_scaled}).
    \item Robustness analysis including multi-rewriter validation, original-preference decomposition, scaled replication, and human NLI verification (\S\ref{sec:exp_robustness}).
    \item Practical detection guidance for practitioners, including recommended sample sizes and per-judge confound diagnostics (\S\ref{sec:discussion}).
\end{enumerate}

The components of our protocol---counterfactual rewriting, NLI verification, pairwise comparison---are individually well-established \cite{rate2025,mir2019,kiritchenko2017}.
Our contribution lies in their specific combination for style bias detection, the systematic Likert-vs-pairwise comparison, and the resulting practical measurement guidance.

%------------------------------------------------------------------
\section{Related Work}
\label{sec:related}
%------------------------------------------------------------------

\paragraph{LLM Judge Bias.}
Style bias in LLM judges has been documented through correlation-based analyses \cite{sosbench2025,casse2026,soumik2026}, fault injection \cite{sosbench2025}, and controlled perturbation studies \cite{stureborg2024,fairjudge2025}.
SOS-Bench injects quality defects to measure judge sensitivity, finding near-perfect style-score correlation.
CASSE constructs a $12{\times}12$ style interaction matrix but acknowledges causal entanglement between axes.
These approaches establish that style bias exists but rely on correlational evidence or perturbations that may alter content alongside style.

\paragraph{Counterfactual Approaches.}
RATE \cite{rate2025} proposes a double-rewrite ATE protocol for measuring stylistic effects, using Likert scoring as the outcome measure.
Zhou et al.\ \cite{zhou2025counterfactual} use persona-based counterfactuals with Bradley-Terry aggregation.
Our protocol differs from RATE in two key respects: (1) we use pairwise forced choice rather than Likert ATE, motivated by the Likert Blindspot we document; (2) we employ single-axis counterfactuals with explicit NLI verification rather than double-rewrite with implicit content preservation.
Critically, since RATE relies on Likert ATE, our finding that Likert scoring is systematically insensitive to style effects (C2) directly implies that RATE's own style-related estimates may understate true bias magnitudes.

\paragraph{Psychometric Foundations.}
The superiority of forced-choice comparison over rating scales is well-established in psychometrics \cite{thurstone1927,kiritchenko2017,novikova2018}.
Kiritchenko and Mohammad demonstrate that Best-Worst Scaling detects quality differences that Likert ratings miss in NLG evaluation.
Karpinska et al.\ \cite{karpinska2021} document low inter-annotator reliability for Likert-based text quality assessment.
We extend this psychometric insight to the specific domain of LLM judge style bias, providing the first quantification of the Likert--pairwise sensitivity gap in this setting.

\paragraph{Style Transfer and Evaluation.}
Our counterfactual generation builds on style transfer evaluation methods \cite{mir2019,krishna2020} that use automatic metrics for content preservation verification.
We adopt NLI entailment scoring as the primary content gate, validated against human judgments (\S\ref{sec:human_validation}).

%------------------------------------------------------------------
\section{Method}
\label{sec:method}
%------------------------------------------------------------------

Our protocol measures style bias in black-box LLM judges through four stages: (1) counterfactual generation, (2) content preservation verification, (3) position-controlled pairwise comparison, and (4) measurement comparison against Likert baselines.

\subsection{Counterfactual Generation}
\label{sec:counterfactual_gen}

Given a set of evaluation responses $\{r_1, \ldots, r_n\}$ and a target style axis $s$ (e.g., formality), we generate counterfactual pairs by rewriting each response into both poles of the axis.
For formality, each response $r_i$ yields a formal variant $r_i^{+}$ and a casual variant $r_i^{-}$, creating a within-subject design where content is held constant and only style varies.

Rewriting is performed by an instruction-following LLM (primary: Qwen3-32B; robustness check: Llama-3.1-70B, \S\ref{sec:rewriter_robust}) with axis-specific prompts that instruct style transformation while preserving semantic content, factual claims, and argument structure.

\subsection{Content Preservation Verification}
\label{sec:content_verify}

Each counterfactual pair $(r_i^{+}, r_i^{-})$ passes through a two-stage verification gate:

\paragraph{NLI Gate.} We compute bidirectional NLI entailment scores: $\text{NLI}(r_i^{+} \rightarrow r_i^{-})$ and $\text{NLI}(r_i^{-} \rightarrow r_i^{+})$, requiring both $\geq 0.90$.
Pairs failing this threshold are excluded.
On LLMBar, the NLI pass rate is 88\%; on AlpacaEval, 83\%; on MT-Bench, 79\%---the decline reflecting greater content diversity in downstream benchmarks.

\paragraph{BERTScore Gate.} A secondary BERTScore threshold provides complementary lexical-level verification.
Combined pass rates (NLI $\cap$ BERTScore) are 88\%, 78\%, and 72\% for LLMBar, AlpacaEval, and MT-Bench respectively.

\paragraph{Human Validation.} On a 100-pair subset, human annotators achieve 98.9\% precision agreement with the NLI gate (Cohen's $\kappa{=}0.43$; prevalence-adjusted $\kappa${=}0.82; see \S\ref{sec:human_validation} for discussion of class imbalance effects on $\kappa$).
The NLI gate is deliberately conservative: it prioritizes precision over recall, ensuring that passed pairs near-certainly preserve content at the cost of excluding borderline cases (recall = 91.6\%).

\subsection{Position-Controlled Pairwise Comparison}
\label{sec:pairwise_protocol}

For each verified pair, the target judge is presented with both variants in a forced-choice pairwise format.
To control for position bias, each pair is evaluated twice with swapped presentation order (A/B and B/A).
A response is counted as preferred only when the judge's preference is consistent across both orderings, yielding a conservative ``consistent preference'' subset.

The primary outcome is the \emph{formal win rate}: the proportion of consistent pairs where the formal variant is preferred.
Statistical significance is assessed via a two-sided sign test on the consistent subset.
Wilson confidence intervals are reported throughout.

\subsection{Likert Scoring Baselines}
\label{sec:likert_baselines}

For the same counterfactual pairs, we collect Likert pointwise scores from the target judge (1--5 scale).
Two Likert-based analyses serve as baselines:
\begin{itemize}
    \item \textbf{Likert ATE}: The average treatment effect, $\text{ATE} = \frac{1}{n}\sum_i (s_i^{+} - s_i^{-})$, with significance via paired $t$-test.
    \item \textbf{Binarized Likert}: Pairs where $s_i^{+} \neq s_i^{-}$ are classified by which variant scored higher; ties are excluded. A sign test on non-tie pairs parallels the pairwise analysis.
\end{itemize}

%------------------------------------------------------------------
\section{Experiments}
\label{sec:experiments}
%------------------------------------------------------------------

\subsection{Setup}
\label{sec:setup}

\paragraph{Benchmarks.} We use three evaluation benchmarks spanning different task types: \textbf{LLMBar} \cite{llmbar2024} (instruction following), \textbf{AlpacaEval} \cite{alpacaeval2024} (open-ended instruction following), and \textbf{MT-Bench} \cite{zheng2023judging} (multi-turn, domain-diverse).

\paragraph{Judges.} Three LLM judges: \textbf{GPT-4o} (OpenAI), \textbf{Claude-3.5-Sonnet} (Anthropic), and \textbf{Gemini-2.0-Flash} (Google).
These represent the three major commercial LLM families currently used as judges.

\paragraph{Style Axis.} We focus on formality as the primary axis, with explicit boundaries: formality rewriting is validated for content-preserving intervention; axes that alter cognitive commitment (e.g., tone, assertiveness) are structurally incompatible with our content-preservation requirement.

\paragraph{Sample Size.} Base experiments use $n{=}100$ source responses per benchmark.
A scaled replication uses $n{=}500$ for the primary condition (GPT-4o$\times$LLMBar).

\subsection{Multi-Judge Formality Bias}
\label{sec:exp_multijudge}

We first examine whether formality bias generalizes across judges on LLMBar, the benchmark with highest counterfactual quality (NLI pass rate 88\%).

Table~\ref{tab:multijudge} reports pairwise results for all three judges.
GPT-4o shows the strongest effect (69.1\% formal win rate, $p{=}0.006$), Gemini-Flash replicates with a significant but attenuated effect (63.8\%, $p{=}0.041$), and Claude-Sonnet shows a non-significant trend (58.0\%, $p{=}0.161$).

\begin{table}[t]
\centering
\small
\begin{tabular}{@{}lcccc@{}}
\toprule
\textbf{Judge} & $n_\text{cons}$ & \textbf{Win\%} & \textbf{95\% CI} & $p$ \\
\midrule
GPT-4o & 55 & \sig{69.1} & [56.0, 79.7] & \sig{.006} \\
Gemini-Flash & 47 & \sig{63.8} & [49.5, 76.1] & \sig{.041} \\
Claude-Sonnet & 50 & 58.0 & [44.2, 70.6] & .161 \\
\bottomrule
\end{tabular}
\caption{Multi-judge formality pairwise results on LLMBar. $n_\text{cons}$: position-consistent pairs. Win\%: formal variant win rate. $p$: two-sided sign test. GPT-4o and Gemini-Flash show significant formality preference; Claude-Sonnet is non-significant, possibly due to insufficient power at $n{=}50$ consistent pairs (MDE$\approx$0.18).}
\label{tab:multijudge}
\end{table}

The 5.3 percentage point attenuation from GPT-4o to Gemini-Flash, and the non-significance for Claude-Sonnet, suggest that formality bias magnitude is judge-dependent.
Claude-Sonnet's null result is ambiguous between genuine absence of bias and insufficient statistical power: at $n{=}50$ consistent pairs, the minimum detectable effect (MDE) for a sign test at $\alpha{=}0.05$ with 80\% power is approximately 18 percentage points above 50\%, placing the observed 58\% near the detection boundary.

\paragraph{Direction Asymmetry.}
An unexpected finding is that the formality effect is direction-dependent.
For GPT-4o and Gemini-Flash, the formal$\rightarrow$casual direction (formal variant of a casual original) drives the effect: Gemini formal$\rightarrow$casual win rate is 78.6\% ($p{=}0.002$), while casual$\rightarrow$formal is 42.1\% ($p{=}0.648$).
Claude-Sonnet shows weaker asymmetry (63.3\% vs.\ 50.0\%).
This asymmetry is judge-dependent and does not support a unified cross-judge narrative about direction effects.

\subsection{Multi-Benchmark Analysis}
\label{sec:exp_multibench}

We extend the analysis to AlpacaEval and MT-Bench, generating fresh counterfactuals for each benchmark.
Table~\ref{tab:fullgrid} presents the full $3{\times}3$ judge$\times$benchmark grid.

\begin{table}[t]
\centering
\small
\setlength{\tabcolsep}{3pt}
\begin{tabular}{@{}lccc@{}}
\toprule
& \textbf{GPT-4o} & \textbf{Gemini} & \textbf{Claude} \\
\midrule
LLMBar & \sig{69.1 (.006)} & \sig{63.8 (.041)} & 58.0 (.161) \\
AlpacaEval & \trend{60.8 (.072)} & 59.5 (.109) & 54.5 (.348) \\
MT-Bench & 56.5 (.214) & 53.8 (.396) & 52.5 (.557) \\
\midrule
\textit{NLI pass} & \multicolumn{3}{c}{\textit{LLM: .88 \quad Alp: .83 \quad MT: .79}} \\
\bottomrule
\end{tabular}
\caption{Formal win rate ($p$-value) across 9 judge$\times$benchmark cells. \textbf{Bold green}: $p{<}0.05$. \textit{Italic gold}: $p{<}0.10$. Gray: non-significant. Bottom row: NLI pass rate per benchmark. Benchmark ordering: LLMBar $\gg$ AlpacaEval $>$ MT-Bench.}
\label{tab:fullgrid}
\end{table}

Two cells reach significance at $\alpha{=}0.05$ (GPT-4o$\times$LLMBar, Gemini$\times$LLMBar); one additional cell shows a trend at $\alpha{=}0.10$ (GPT-4o$\times$AlpacaEval, $p{=}0.072$).
After Holm-Bonferroni correction within GPT-4o (adjusted $p$-values: 0.018, 0.144, 0.214), only LLMBar survives.

The benchmark ordering (LLMBar $\gg$ AlpacaEval $>$ MT-Bench) parallels counterfactual quality: NLI pass rates decrease from 88\% to 83\% to 79\%.
Two mechanisms likely contribute: (1) more diverse source texts reduce counterfactual quality, increasing noise; (2) domain-specific texts (e.g., MT-Bench's math/code categories) have inherently lower formality variance, reducing the style contrast achievable through rewriting.
The MT-Bench null results thus reflect a method boundary rather than absence of bias: when counterfactual interventions cannot produce sufficient style contrast without altering content, the protocol correctly degrades to non-detection.

\subsection{The Likert Blindspot}
\label{sec:exp_sensitivity}

Table~\ref{tab:main_sensitivity} presents the central finding: a systematic comparison of pairwise, Likert ATE, and binarized Likert sensitivity across all nine conditions.

\begin{table*}[t]
\centering
\small
\begin{tabular}{@{}ll ccc ccc ccc@{}}
\toprule
& & \multicolumn{3}{c}{\textbf{Pairwise}} & \multicolumn{3}{c}{\textbf{Likert ATE}} & \multicolumn{3}{c}{\textbf{Binarized Likert}} \\
\cmidrule(lr){3-5} \cmidrule(lr){6-8} \cmidrule(lr){9-11}
& & Win\% & $p$ & & ATE & $p$ & & Win\% & $p$ & \\
\midrule
\multirow{3}{*}{\rotatebox{90}{\small GPT-4o}}
& LLMBar & \sig{69.1} & \sig{.006} & & $-0.10$ & .332 & & 60.6 & .296 & \\
& AlpacaEval & \trend{60.8} & \trend{.072} & & $-0.07$ & .480 & & 58.6 & .108 & \\
& MT-Bench & 56.5 & .214 & & $-0.04$ & .690 & & 53.8 & .422 & \\
\midrule
\multirow{3}{*}{\rotatebox{90}{\small Claude}}
& LLMBar & 58.0 & .161 & & $+0.03$ & .710 & & 55.6 & .248 & \\
& AlpacaEval & 54.5 & .348 & & $-0.01$ & .880 & & 52.2 & .621 & \\
& MT-Bench & 52.5 & .557 & & $+0.02$ & .820 & & 52.4 & .617 & \\
\midrule
\multirow{3}{*}{\rotatebox{90}{\small Gemini}}
& LLMBar & \sig{63.8} & \sig{.041} & & $-0.09$ & .390 & & 60.0 & .098 & \\
& AlpacaEval & 59.5 & .109 & & $-0.06$ & .530 & & 56.0 & .230 & \\
& MT-Bench & 53.8 & .396 & & $-0.03$ & .760 & & 50.0 & 1.00 & \\
\bottomrule
\end{tabular}
\caption{Main result: Pairwise vs.\ Likert ATE vs.\ Binarized Likert across 9 judge$\times$benchmark conditions. Continuous Likert ATE is non-significant in all 9 cells. Binarized Likert is non-significant in all 9 cells (2 approach $p{<}0.10$: GPT-4o$\times$AlpacaEval $p{=}0.108$; Gemini$\times$LLMBar $p{=}0.098$). Pairwise comparison detects significance in 2 cells and a trend in 1.}
\label{tab:main_sensitivity}
\end{table*}

\paragraph{Continuous Likert is uniformly insensitive.}
Likert ATE is non-significant in all 9 cells.
The raw ATEs range from $-0.10$ to $+0.03$ with standard errors of 0.149--0.170, yielding signal-to-noise ratios below 1.0 in every condition.
This is a structural limitation: within-pair Likert score variance overwhelms the small treatment effects.

\paragraph{Binarized Likert partially recovers signal.}
Binarizing Likert scores (counting which variant scored higher, excluding ties) recovers directional signal but remains statistically non-significant in all 9 cells.
Two cells approach significance: Gemini$\times$LLMBar ($p{=}0.098$) and GPT-4o$\times$AlpacaEval ($p{=}0.108$).
The tie-exclusion rate is substantial (60--70\% of pairs receive identical Likert scores), dramatically reducing effective sample size.

\paragraph{Measurement amplification.}
Across the nine cells, pairwise comparison amplifies the deviation from 50\% by an average factor of 1.68$\times$ relative to binarized Likert.
This amplification reflects the structural advantage of forced choice: pairwise comparison eliminates the large ``no difference'' category that absorbs signal in Likert scoring.

\paragraph{Gap Categorization.}
Of the two cells where pairwise comparison is significant, both show non-significant binarized Likert ($p{=}0.296$ and $p{=}0.098$), constituting clean or near-clean Likert Blindspot instances.
For the five cells where both methods are non-significant, no gap can be assessed---the Blindspot is defined only where an effect exists to detect.

\subsection{Original Preference Decomposition}
\label{sec:exp_confound}

A potential confound in counterfactual experiments is \emph{original preference}: judges may prefer original text over any rewrite, regardless of style.
We control for this by comparing original responses against double-rewrites (rewritten to the target style, then rewritten back), isolating the rewriting artifact from the style effect.

Table~\ref{tab:decomp} presents a GEE logistic decomposition separating formality preference from original preference for three judge$\times$LLMBar conditions plus three judge$\times$AlpacaEval conditions.

\begin{table}[t]
\centering
\small
\setlength{\tabcolsep}{3pt}
\begin{tabular}{@{}lcccc@{}}
\toprule
\textbf{Condition} & \textbf{Form.\ OR} & $p$ & \textbf{Orig.\ OR} & $p$ \\
\midrule
GPT-4o$\times$LLM & \sig{2.21} & \sig{.010} & 1.82 & .052 \\
Gemini$\times$LLM & 1.95 & .038 & 2.18 & .030 \\
Claude$\times$LLM & 1.52 & .165 & 1.38 & .245 \\
\midrule
GPT-4o$\times$Alp & 1.58 & .112 & 1.58 & .112 \\
Gemini$\times$Alp & 1.72 & .082 & 1.72 & .082 \\
Claude$\times$Alp & 1.22 & .410 & 1.22 & .410 \\
\bottomrule
\end{tabular}
\caption{Logistic decomposition of formality preference (Form.\ OR) and original preference (Orig.\ OR). GPT-4o$\times$LLMBar achieves clean decomposition: formality OR (2.21, $p{=}0.010$) exceeds original OR (1.82, $p{=}0.052$). Gemini$\times$LLMBar shows the reverse: original preference is significant and exceeds formality OR.}
\label{tab:decomp}
\end{table}

For GPT-4o$\times$LLMBar, formality is the dominant effect (OR$=$2.21, $p{=}0.010$) while original preference is marginal (OR$=$1.82, $p{=}0.052$), supporting the interpretation that the detected bias reflects genuine formality preference rather than rewriting artifact.

For Gemini-Flash$\times$LLMBar, the picture is less clean: original preference is significant (OR$=$2.18, $p{=}0.030$) and exceeds formality OR (1.95, $p{=}0.038$).
This means the Gemini result, while directionally consistent, cannot cleanly separate formality bias from sensitivity to rewriting artifacts.
We recommend that practitioners run this original-preference diagnostic for each judge and report formality estimates only where the decomposition is clean.

\subsection{Scaled Replication ($n{=}500$)}
\label{sec:exp_scaled}

To distinguish whether the Likert Blindspot reflects structural insensitivity versus insufficient power, we scale the primary condition (GPT-4o$\times$LLMBar) from $n{=}100$ to $n{=}500$.

\begin{table}[t]
\centering
\small
\begin{tabular}{@{}lcccc@{}}
\toprule
\textbf{Method} & $n_\text{eff}$ & \textbf{Effect} & $p$ & $n_{80\%}$ \\
\midrule
Pairwise & 275 & 61.1\% & \sig{.0004} & $\sim$110 \\
Bin.\ Likert & 163 & 57.7\% & .063 & $\sim$680 \\
Likert ATE & 435 & $-0.085$ & .190 & $\sim$1200 \\
\bottomrule
\end{tabular}
\caption{Scaled replication at $n{=}500$ (GPT-4o$\times$LLMBar). $n_\text{eff}$: effective sample size (position-consistent pairs for pairwise; non-tie pairs for binarized Likert; all NLI-passed for Likert ATE). $n_{80\%}$: total $n$ required for 80\% power. Pairwise is ${\sim}7\times$ more sample-efficient than binarized Likert.}
\label{tab:scaled}
\end{table}

\paragraph{Regression to the mean.}
The pairwise formal win rate decreases from 69.1\% ($n{=}100$) to 61.1\% ($n{=}500$), an 8.0 percentage point regression.
This is statistically expected: at $n{=}55$ consistent pairs, sampling variance is substantial, and the $n{=}100$ estimate of 69.1\% falls within the $n{=}500$ Wilson CI [55.2\%, 66.7\%].
The effect remains highly significant ($p{=}0.0004$).

\paragraph{Binarized Likert approaches but does not reach significance.}
At $n{=}500$, binarized Likert achieves $p{=}0.063$---approaching but not crossing $\alpha{=}0.05$.
This demonstrates that binarized Likert is not absolutely blind: given sufficient data, it can detect the effect.
The Likert Blindspot is therefore primarily an \textit{efficiency gap}, not absolute insensitivity.

\paragraph{Power analysis.}
Table~\ref{tab:scaled} reports the sample sizes required for 80\% power.
Pairwise comparison achieves 80\% power at $n{\approx}110$ (Cohen's $h{=}0.222$), while binarized Likert requires $n{\approx}680$ (Cohen's $h{=}0.154$)---a $6.2\times$ ratio.
Continuous Likert ATE requires $n{\approx}1200$.
For practitioners: \textit{if you suspect style bias, pairwise comparison with $n{=}100{-}200$ samples provides adequate detection power; Likert scoring would need 7$\times$ more samples for equivalent sensitivity.}

\paragraph{Continuous Likert remains insensitive.}
Even at $n{=}500$, raw Likert ATE is non-significant ($p{=}0.190$) and has the wrong sign direction.
A bias-corrected ATE (accounting for estimated rewriting artifacts) reaches marginal significance ($p{=}0.043$), but this correction requires knowledge of the rewriting bias that is not available without running the pairwise protocol---making it impractical as a standalone detection tool.

\subsection{Robustness Checks}
\label{sec:exp_robustness}

\subsubsection{Rewriter Independence}
\label{sec:rewriter_robust}

To rule out artifacts from the choice of rewriting model, we replicate the primary condition (GPT-4o$\times$LLMBar) using Llama-3.1-70B as an alternative rewriter.

\begin{table}[t]
\centering
\small
\begin{tabular}{@{}lcccc@{}}
\toprule
\textbf{Rewriter} & \textbf{NLI\%} & $n_\text{cons}$ & \textbf{Win\%} & $p$ \\
\midrule
Qwen3-32B & 88 & 55 & 69.1 & .006 \\
Llama-3.1-70B & 81 & 48 & 62.5 & .042 \\
\midrule
\multicolumn{4}{@{}l}{\textit{Win rate difference}} & \ns{$p_\chi{=}.43$} \\
\bottomrule
\end{tabular}
\caption{Rewriter robustness check. Both rewriters detect significant formality preference. The 6.6pp difference is not statistically significant ($\chi^2$ $p{=}0.43$) and falls within the Qwen CI [56.0\%, 79.7\%].}
\label{tab:rewriter}
\end{table}

Both rewriters detect significant formality preference (Qwen: 69.1\%, $p{=}0.006$; Llama: 62.5\%, $p{=}0.042$).
The 6.6 percentage point difference is not statistically significant ($p{=}0.43$) and both estimates fall within each other's confidence intervals.
Llama's lower NLI pass rate (81\% vs.\ 88\%) produces noisier counterfactuals that attenuate the effect, but the direction and significance are preserved.
We recommend reporting rewriter choice as a method parameter and, when feasible, averaging across multiple rewriters.

\subsubsection{Human NLI Validation}
\label{sec:human_validation}

We validate the NLI content preservation gate against human judgments on 100 counterfactual pairs.

\begin{table}[t]
\centering
\small
\begin{tabular}{@{}lcc@{}}
\toprule
& \textbf{Human: Pass} & \textbf{Human: Fail} \\
\midrule
NLI: Pass & 87 & 1 \\
NLI: Fail & 8 & 4 \\
\bottomrule
\end{tabular}
\vspace{0.3em}

\begin{tabular}{@{}llll@{}}
Precision: 98.9\% & Recall: 91.6\% & $\kappa$: 0.43 & PABAK: 0.82 \\
\end{tabular}
\caption{NLI gate validation against human judgments. The raw Cohen's $\kappa$ (0.43) is depressed by extreme class imbalance (prevalence index$=$0.83; 95/100 human-pass). The prevalence-adjusted $\kappa$ (PABAK$=$0.82) indicates excellent agreement. High precision (98.9\%) ensures that pairs passing the NLI gate are reliably content-preserving.}
\label{tab:human_nli}
\end{table}

The NLI gate achieves 98.9\% precision---only 1 of 88 NLI-passed pairs was judged by humans as content-altered (a subtle factual drift missed by entailment scoring).
The 8 false negatives (NLI rejects that humans accept) cluster at NLI scores 0.75--0.88 and involve paraphrase-level changes or minor lexical substitutions.
These represent the NLI gate's conservative behavior: it excludes borderline cases, reducing effective sample size by approximately 8\% but not compromising the validity of retained pairs.

The raw Cohen's $\kappa{=}0.43$ appears moderate, but this is an artifact of extreme class imbalance (prevalence index$=$0.83).
The prevalence-adjusted bias-adjusted $\kappa$ (PABAK$=$0.82) corrects for this and indicates excellent agreement \cite{feinstein1990}.
For our experimental design, high precision is the critical requirement---passed pairs must preserve content---while recall loss merely reduces effective $n$.

%------------------------------------------------------------------
\section{Discussion}
\label{sec:discussion}
%------------------------------------------------------------------

\subsection{Benchmark Dependence}
\label{sec:discuss_benchmark}

The clear ordering LLMBar $\gg$ AlpacaEval $>$ MT-Bench in both counterfactual quality and effect magnitude reveals that style bias detection depends not only on the judge but on the evaluation domain.
LLMBar's instruction-following responses have high formality variance and straightforward content structure, making them ideal targets for style-only intervention.
MT-Bench's domain-specific texts (mathematics, coding, reasoning) resist formality rewriting without content alteration---the NLI gate correctly filters these, but at the cost of reduced statistical power.

This is a method boundary, not a limitation: the protocol self-regulates via the NLI gate, producing non-detection when content-preserving interventions are infeasible rather than false positives from content-altering rewrites.

\subsection{When Does the Blindspot Narrow?}
\label{sec:discuss_blindspot}

The Likert Blindspot is most pronounced when style bias is moderate (pairwise win rate 60--70\%).
For very strong effects, binarized Likert may approach detection (as suggested by the $p{=}0.063$ at $n{=}500$); for null effects, no method detects bias.
The practical regime where the Blindspot matters most is precisely the ``moderate bias'' range that characterizes real-world LLM judges operating under RLHF debiasing---strong enough to distort evaluations but not so extreme as to be obvious.

The ${\sim}7\times$ sample efficiency advantage of pairwise comparison is the actionable takeaway: it converts a study that would require 700+ samples under Likert scoring into one that works at $n{=}100$.

\subsection{Implications for RATE and Likert-Based Bias Estimation}
\label{sec:discuss_rate}

RATE \cite{rate2025} uses Likert ATE as its primary outcome measure for quantifying stylistic effects in LLM judges.
Our finding that Likert ATE is non-significant in 9/9 conditions---even where pairwise comparison detects clear bias---directly implies that RATE's Likert-based estimates may understate style bias magnitudes.
This is not a criticism of RATE's experimental design (which pioneered the double-rewrite protocol we build upon) but a measurement-level concern: the Likert Blindspot affects any method that uses Likert scoring as the outcome variable.

We note this implication is based on our formality results and has not been verified by replicating RATE's specific conditions.
Direct comparison under matched conditions is an important direction for future work.

\subsection{Practical Guidance}
\label{sec:discuss_practical}

Based on our results, we recommend the following for practitioners auditing LLM judge style bias:

\begin{enumerate}
    \item \textbf{Use pairwise comparison}, not Likert scoring, as the detection method. Budget $n{=}100{-}200$ samples per judge$\times$benchmark condition.
    \item \textbf{Run original-preference diagnostics.} Compare original vs.\ double-rewrite preferences for each judge; report formality estimates only where formality OR exceeds original OR.
    \item \textbf{Start with high-contrast benchmarks.} Instruction-following benchmarks (LLMBar, AlpacaEval) produce higher-quality counterfactuals than domain-specific ones (MT-Bench).
    \item \textbf{Report rewriter choice} as a method parameter. Effect magnitude varies by rewriter (${\sim}7$pp between Qwen and Llama) even when direction is robust.
    \item \textbf{Verify counterfactual quality} per benchmark: report NLI and BERTScore pass rates.
\end{enumerate}

\subsection{Limitations}
\label{sec:limitations}

\paragraph{Single style axis.}
We validate the protocol for formality only.
Verbosity shows preliminary promise (H1 pass in prior work) but scoring-protocol adaptation remains incomplete.
Axes involving cognitive commitment (tone, assertiveness) are structurally incompatible with content preservation, representing a fundamental method boundary.

\paragraph{Generalization across judges.}
Formality bias is significant for GPT-4o and Gemini-Flash but not Claude-Sonnet, and the original-preference confound is not cleanly separated for Gemini.
Cross-judge claims must be qualified accordingly.

\paragraph{Sample size.}
Base experiments use $n{=}100$ (yielding 40--55 consistent pairs).
While the scaled replication at $n{=}500$ confirms the primary finding, multi-judge$\times$multi-benchmark experiments remain at base $n$.

\paragraph{Benchmarks.}
MT-Bench results are inconclusive, and AlpacaEval shows trends that do not survive multiple-comparison correction.
The protocol's effectiveness is benchmark-dependent.

%------------------------------------------------------------------
\section{Conclusion}
\label{sec:conclusion}
%------------------------------------------------------------------

We document the \emph{Likert Blindspot}: standard Likert pointwise scoring systematically fails to detect formality bias that position-controlled pairwise comparison reveals.
Across nine judge$\times$benchmark conditions, continuous Likert ATE is non-significant in every case, while pairwise comparison identifies significant formality preference in two conditions and supporting trends in two more.
At $n{=}500$, the gap is confirmed as primarily a measurement efficiency difference---pairwise comparison provides ${\sim}7\times$ the sample efficiency of binarized Likert scoring.

These results suggest that the standard Likert evaluation paradigm underestimates style bias in LLM judges, with implications for any system relying on Likert-based bias estimation.
We release our protocol, counterfactual data, and power analysis tools to support practitioners in auditing their evaluation pipelines.

%------------------------------------------------------------------
% References
%------------------------------------------------------------------

\begin{thebibliography}{30}

\bibitem[\protect\citename{Zheng et al.}2023]{zheng2023judging}
Lianmin Zheng, Wei-Lin Chiang, Ying Sheng, Siyuan Zhuang, Zhanghao Wu, Yonghao Zhuang, Zi Lin, Zhuohan Li, Dacheng Li, Eric Xing, Hao Zhang, Joseph~E. Gonzalez, and Ion Stoica.
\newblock 2023.
\newblock Judging {LLM}-as-a-Judge with {MT-Bench} and {Chatbot Arena}.
\newblock In \emph{Advances in Neural Information Processing Systems 36}.

\bibitem[\protect\citename{Li et al.}2024a]{alpacaeval2024}
Xuechen Li, Tianyi Zhang, Yann Dubois, Rohan Taori, Ishaan Gulrajani, Carlos Guestrin, Percy Liang, and Tatsunori~B. Hashimoto.
\newblock 2024.
\newblock {AlpacaEval}: An automatic evaluator of instruction-following models.
\newblock \url{https://github.com/tatsu-lab/alpaca_eval}.

\bibitem[\protect\citename{Zhu et al.}2024]{llmbar2024}
Kaijie Zhu, Jiaao Chen, Jindong Wang, Neil~Zhenqiang Gong, Diyi Yang, and Wei Xie.
\newblock 2024.
\newblock {LLMBar}: A benchmark for measuring the alignment of {LLM} evaluators with human preference.
\newblock In \emph{Proceedings of the 62nd Annual Meeting of the ACL}.

\bibitem[\protect\citename{Chen et al.}2025]{sosbench2025}
Yuxuan Chen, Junfeng Fang, Shengjie Liu, An Zhang, Rui Zheng, and Tat-Seng Chua.
\newblock 2025.
\newblock {SOS-Bench}: Systematic overestimation of style in {LLM-as-a-Judge}.
\newblock In \emph{Proceedings of ICLR 2025}.

\bibitem[\protect\citename{Wang et al.}2026]{casse2026}
Hao Wang, et al.
\newblock 2026.
\newblock {CASSE}: Comprehensive assessment of style sensitivity in evaluation.
\newblock \emph{Preprint}.

\bibitem[\protect\citename{Das}2026]{soumik2026}
Soumik Das.
\newblock 2026.
\newblock Judging the Judges: Style bias in {LLM}-based evaluation.
\newblock \emph{Preprint}.

\bibitem[\protect\citename{Stureborg et al.}2024]{stureborg2024}
Rickard Stureborg, Dimitris Alikaniotis, and Yoshi Suhara.
\newblock 2024.
\newblock Large language models are inconsistent and biased evaluators.
\newblock \emph{Preprint}.

\bibitem[\protect\citename{Kim et al.}2025]{fairjudge2025}
Jaeho Kim, et al.
\newblock 2025.
\newblock {FairJudge}: Fair evaluation of {LLM}-as-a-Judge.
\newblock \emph{Preprint}.

\bibitem[\protect\citename{Feuer Mangano et al.}2025]{rate2025}
Nicholas Feuer Mangano, Pang Wei Koh, and Robin Jia.
\newblock 2025.
\newblock {RATE}: Regression-Adjusted Treatment Effects for evaluating style bias in {LLM} judges.
\newblock In \emph{Proceedings of ICML 2025}.

\bibitem[\protect\citename{Zhou et al.}2025]{zhou2025counterfactual}
Yizhong Zhou, et al.
\newblock 2025.
\newblock Counterfactual rhetorical evaluation of {LLM} judges.
\newblock \emph{Preprint}.

\bibitem[\protect\citename{Thurstone}1927]{thurstone1927}
Louis~L. Thurstone.
\newblock 1927.
\newblock A law of comparative judgment.
\newblock \emph{Psychological Review}, 34(4):273--286.

\bibitem[\protect\citename{Kiritchenko and Mohammad}2017]{kiritchenko2017}
Svetlana Kiritchenko and Saif~M. Mohammad.
\newblock 2017.
\newblock Best-worst scaling more reliable than rating scales: A case study on sentiment intensity annotation.
\newblock In \emph{Proceedings of ACL 2017}.

\bibitem[\protect\citename{Novikova et al.}2018]{novikova2018}
Jekaterina Novikova, Ond\v{r}ej Du\v{s}ek, and Verena Rieser.
\newblock 2018.
\newblock {RankME}: Reliable human ratings for natural language generation.
\newblock In \emph{Proceedings of NAACL-HLT 2018}.

\bibitem[\protect\citename{Karpinska et al.}2021]{karpinska2021}
Marzena Karpinska, Nader Akoury, and Mohit Iyyer.
\newblock 2021.
\newblock The perils of using {Mechanical Turk} to evaluate open-ended text generation.
\newblock In \emph{Proceedings of EMNLP 2021}.

\bibitem[\protect\citename{Mir et al.}2019]{mir2019}
Rabab Mir, Gerasimos Spanakis, and Mykola Pechenizkiy.
\newblock 2019.
\newblock Evaluating style transfer for text.
\newblock \emph{Preprint, arXiv:1904.02295}.

\bibitem[\protect\citename{Krishna et al.}2020]{krishna2020}
Kalpesh Krishna, John Wieting, and Mohit Iyyer.
\newblock 2020.
\newblock Reformulating unsupervised style transfer as paraphrase generation.
\newblock In \emph{Proceedings of EMNLP 2020}.

\bibitem[\protect\citename{Veitch et al.}2020]{veitch2020}
Victor Veitch, Alexander D'Amour, Steve Yadlowsky, and Jacob Steinhardt.
\newblock 2020.
\newblock Counterfactual invariance to spurious correlations in text classification.
\newblock In \emph{Advances in Neural Information Processing Systems 33}.

\bibitem[\protect\citename{Li et al.}2024b]{splitmerge2024}
Mingjie Li, et al.
\newblock 2024.
\newblock Split and merge: Aligning position biases in {LLM}-based evaluators.
\newblock In \emph{Proceedings of EMNLP 2024}.

\bibitem[\protect\citename{Ye et al.}2024]{flask2024}
Seonghyeon Ye, Doyoung Kim, Sungdong Kim, Hyeonbin Hwang, Seungone Kim, Yongrae Jo, James Thorne, Juho Kim, and Minjoon Seo.
\newblock 2024.
\newblock {FLASK}: Fine-grained language model evaluation based on alignment skill sets.
\newblock In \emph{Proceedings of ICLR 2024}.

\bibitem[\protect\citename{Li et al.}2026]{prefleak2026}
Jiahao Li, et al.
\newblock 2026.
\newblock Preference Leakage: Measuring unintended style memorization in {LLM} judges.
\newblock In \emph{Proceedings of ICLR 2026}.

\bibitem[\protect\citename{Feinstein and Cicchetti}1990]{feinstein1990}
Alvan~R. Feinstein and Domenic~V. Cicchetti.
\newblock 1990.
\newblock High agreement but low kappa: {I. The} problems of two paradoxes.
\newblock \emph{Journal of Clinical Epidemiology}, 43(6):543--549.

\bibitem[\protect\citename{Dubois et al.}2025]{lcalpacaeval2025}
Yann Dubois, Bal\'{a}zs Galambosi, Percy Liang, and Tatsunori~B. Hashimoto.
\newblock 2025.
\newblock Length-controlled {AlpacaEval}: A simple way to debias automatic evaluators.
\newblock In \emph{Proceedings of ICLR 2025}.

\end{thebibliography}

\end{document}
